\scanpage{72}
\mspage{- \handwritten{6} -}\hfill n\textsuperscript{o} \handwritten{380}

3. \uline{Inverses}

Dans ce N\textsuperscript{o} on suppose que l'algèbre $A$ est normée \uline{complète} et possède un élément unité $e$.

PROPOSITION 4.- \uline{Pour tout élément} $x\handwritten{\in}A$ \uline{tel que} $\handwritten{\unsure{\mathscr{D}|\mathfrak{D}|\partial}}(e-x)<1$, \uline{la série} $\displaystyle\handwritten{\sum}_{u=1}^{\handwritten{\infty}}(e-x)^n$ \uline{est absolument convergente et} $\displaystyle e+\handwritten{\sum}_{u=1}^{\handwritten{\infty}}(e-x)^n$ \uline{est l'inverse de} $x$.

La série est absolument convergente d'après le Critère de Cauchy. Soit $\displaystyle x'=e+\handwritten{\sum}_{u=1}^{\handwritten{\infty}}(e-x)^n$. On a
\[
x'x=xx'=x+\handwritten{\sum}_{u=1}^{\handwritten{\infty}}(e-x)^n x=x+\handwritten{\sum}_{u=1}^{\handwritten{\infty}}(e-x)^n-\handwritten{\sum}_{u=1}^{\handwritten{\infty}}(e-x)^{n+1}=e.
\]
COROLLAIRE.- \uline{Le groupe des éléments inversibles de} $A$ \uline{contient la boule ouverte de centre} $e$ \uline{de rayon} $1$.

En effet $\handwritten{\mathscr{D}}(x-e)\handwritten{\le}\lVert x-e\rVert$ pour tout $x$.

Ce résultat figurait en Top. Gén., Chap.9, \handwritten{\S} 3, prop.13. On a vu par la même occasion que le groupe des éléments inversibles de $A$ est un ouvert de $A$ et que la topologie induite sur ce groupe par la topologie de $A$ est compatible avec la structure de groupe. Si $A$ est de dimension finie sur $C$, ce groupe ne possède qu'un nombre fini de composantes connexes. Dans les autres cas, il arrive qu'il possède un nombre infini de composantes connexes.

PROPOSITION 5.- \uline{Soit} $A$ \unsure{'}\uline{une algèbre normée complète avec unité} $e$ \uline{et soit} $x_n$ \uline{une suite convergente d'éléments inversibles de} $A$ \uline{ayant pour limite} $x$ :

a) \uline{si la suite} $\lVert x_n^{-1}\rVert$ \uline{est bornée}, \uline{alors} $x$ \uline{est inversible},

a') \uline{si la suite} $\handwritten{\mathscr{D}}(x_n^{-1})$ \uline{est bornée et si} $xx_n=x_nx$ \uline{pour tout} $n$, \uline{alors} $x$ \uline{est inversible}.

\uline{Dans les deux cas}, \uline{l'inverse de} $x$ \uline{est la limite des} $x_n^{-1}$.
% archive sheet number: 71

